\section{Carrier Refinement}
\label{sec:refinement}

\subsection{Problems of the naive inversion}
\label{sec:naive-problems}

The fused Gaussians of \cref{sec:tomography} have accurate means early, but their
remaining parameters are only geometrically motivated.  Handing them directly to a
standard optimizer with densification makes the optimizer replace the initialization
instead of refining it, and the experiment then measures the recovery ability of
densification and not the initializer.  Three concrete mismatches cause this.  First, the
fused covariance targets the raw tangent lift, while the renderer displays
$P_{iv} = J_{iv} \Sigma_i J_{iv}\T + 0.3\,I_2$ with the EWA dilation.  The initialization
is thus inconsistent with what the optimizer will render.  Second, opacity cannot be
inferred from the captures, since a fitted 2D amplitude is a normalized-blend gauge
quantity, see \cref{sec:gauge}, and does not identify a per-Gaussian 3D alpha under
compositing.  Third, false contributor associations and silhouette leakage place a small
fraction of components outside the fitted image support.

We therefore treat the fused Gaussians as carriers and mature them with a short,
compact-only schedule before any comparison.  The accepted sequence consists of five
stages.\evidence{docs/ADR-002-carrier-refinement.md, ara/logic/claims.md C30, C31}
\begin{enumerate}[leftmargin=1.6em,itemsep=1pt]
  \item Fit-only Beam Fusion from the reconstruction inputs with $K \ge 3$ and no free
        primitive birth.
  \item Renderer-aware covariance repair, see \cref{sec:cov-repair}.
  \item 380 compact fixed-topology SH0 updates with means, rotations, scales, opacity and
        SH0 all trainable.
  \item Strict projected-center pruning against every fitting view, see
        \cref{sec:containment}.
  \item 380 further compact SH0 updates with means frozen, so the pruning invariant cannot
        be undone.
\end{enumerate}
No clone, split, insertion, densification, opacity reset, higher-SH phase or dense
handover exists on this path, and its entry point accepts no image-bearing argument.

\subsection{Covariance repair}
\label{sec:cov-repair}

For carrier $i$ with contributor set $\mathcal{V}_i$ from the lineage and observed 2D
covariances $C_{iv}$ we define the whitened projection
\begin{equation}
  M_{iv} \;=\; C_{iv}^{-1/2}\,\big(J_{iv}\, \Sigma_i\, J_{iv}\T + 0.3\, I_2\big)\,
               C_{iv}^{-1/2}
  \label{eq:whitened}
\end{equation}
and minimize, over $\Sigma_i$ only, the symmetric generalized-log residual
\begin{equation}
  r_i \;=\; \sqrt{\;\frac{1}{|\mathcal{V}_i|}\sum_{v \in \mathcal{V}_i}
        \operatorname{mean}_k \log^2 \lambda_k(M_{iv})\;}
  \label{eq:cov-residual}
\end{equation}
with a Huber data term, a prior toward the beam covariance and explicit SPD and aspect
bounds.  Means, opacity, appearance, lineage and count stay bit-frozen.  Two properties
of \cref{eq:whitened,eq:cov-residual} matter.  First, the target includes the renderer's
$0.3$-pixel dilation, so the repair aims at what will actually be displayed.  Some fitted
targets are sharper than this floor, so zero residual is not always attainable and the
repair aims at consistency and not perfection.  Second, the log-eigenvalue form penalizes
a projection twice as large and half as large equally.  The obvious alternative, a
one-sided whitened Frobenius residual, penalizes expansion more than collapse and omits
the dilation.  It is shrink-biased and cannot serve as the repair objective here.  The
ablation study quantifies the difference, see \cref{sec:ablations}.  On the development
scene the repair of \cref{eq:cov-residual} lowers the validation risk $\Jq$ by $63.2\%$
against optimizing the raw fused output directly, with three of three paired-root
wins.\evidence{ara/logic/claims.md C30}  A repair stage must justify itself by its effect
on the downstream risk and never by its own residual.

Two further repairs are deliberately absent.  Opacity repair is invalid due to the gauge
argument above, and cross-view amplitude-weighted appearance repair shares the same gauge
problem.  The SH0 colors are instead co-adapted by the compact phases below.  The carrier
path also grows no topology.  Even a clone which preserves optical density, using
$a' = 1 - \sqrt{1-a}$ where a plain copy would change coincident alpha from $a$ to
$1-(1-a)^2$, did not pass its ablation gate.\evidence{ara/logic/claims.md C30}

\subsection{Refinement schedule}
\label{sec:phases}

Both compact phases minimize the sampled objective of \cref{eq:sampled-loss} at fixed
topology and SH degree 0.  Phase one trains all parameter families jointly.  Due to the
alpha and color gauge freedom of the compositor, opacity and support cannot be identified
independently, so scale, orientation, opacity and color must co-adapt.  Optimizing means
alone reaches $4.18$ times the all-family phase-one $\Jq$ and is rejected.  A second
all-family phase lowers $\Jq$ to $0.717$ times its stopping value with three of three
paired wins.  Freezing only the means in phase two costs $1.47\%$ $\Jq$ and passes the
frozen $2\%$ non-inferiority gate, which is exactly what makes the containment invariant
of \cref{sec:containment} permanent.  Incremental SH3 improves $\Jq$ by only $2.3$ to
$3.0\%$, below the $5\%$ materiality floor, so the carrier stays at SH0.  All of these
comparisons appear as rows of the ablation study.\evidence{ara/logic/claims.md C30}

\subsection{Silhouette containment}
\label{sec:containment}

Raw masks are forbidden inside the compact paths.  Their legal surrogate is the union of
fitted 2D Gaussian support ellipses.  After phase one, carrier $i$ is retained only if for
every fitting view $v$
\begin{equation}
  z_{iv} > z_{\text{near}}
  \qquad\text{and}\qquad
  \min_j \big(\cam_v(\bm{\mu}_i) - \bm{m}_{vj}\big)\T C_{vj}^{-1}
         \big(\cam_v(\bm{\mu}_i) - \bm{m}_{vj}\big) \;\le\; 3^2,
  \label{eq:containment}
\end{equation}
i.e.\ the projected center must fall inside some positive-amplitude fitted component's
$3\sigma$ ellipse in every view.  The threshold is fixed and cannot be weakened through
configuration.  Phase two freezes the means and rechecks \cref{eq:containment}, so the
invariant survives to the output.  On the development scene the prune removes $5.3$ to
$5.4\%$ of the carriers, retains $4{,}729$ to $4{,}734$, ends with exactly zero violations
and costs $2.69\%$ $\Jq$ against the unpruned anchor.\evidence{ara/logic/claims.md C30, C31}

\Cref{eq:containment} is a visual-hull center invariant.  It removes carriers whose
centers leave the fitted silhouette cone.  It cannot detect a floater inside the
multi-view hull and it does not prove surface occupancy.  Containing the whole Gaussian in
the mask is impossible in principle, since Gaussian support is infinite, and a
finite-footprint rule would erode legitimate primitives at silhouette boundaries.  A
count-preserving soft variant is predeclared for the ablation grid, namely the barrier
\begin{equation}
  L_{\text{support}} = \sum_{i,v}
    \softplus\!\Big(\tfrac{\min_j q_{ivj} - 9}{\tau}\Big)^{2}
    + \lambda_z\, \softplus\!\Big(\tfrac{z_{\text{near}} - z_{iv}}{\tau_z}\Big)^{2},
  \label{eq:soft-support}
\end{equation}
optionally with deterministic samples on the projected $3\sigma$ ellipse to constrain the
footprint instead of only the center.
\todo{\Cref{eq:soft-support} is a design to freeze before use and not an evaluated
method.  An independently measurable surface-occupancy test for hull-interior floaters
remains open, see T-12, so the absence of physical free-floaters must not be claimed from
\cref{eq:containment} alone.}

\begin{figure}[t]
  \centering
  \figslot{5.5cm}{%
    \textbf{Figure 8, refinement stages.}  One row per stage (raw fusion, after
    covariance repair, after phase one, after prune, after phase two).  \textbf{Left:}
    rendered front view.  \textbf{Right:} side view.  Annotate carrier count and $\Jq$
    per stage.  The covariance-repair row should visibly de-elongate the beams.  Uses the
    representative-root artifacts of the 2026-07-28 compact-only runs under
    \repopath{runs/}, re-rendered with the repository viewer.}
  \caption{\redcaption{Stage-by-stage carrier refinement renders.  Producible from
    existing development artifacts.}}
  \label{fig:carrier-stages}
\end{figure}
